% !PRINTED none
\begin{center}
\gesperrt{Twentieth Chapter.}

\gesperrt{Quadratic Fields.}
\end{center}

\begin{center}
\rule{0.2\linewidth}{0.4pt}
\end{center}

\begin{center}
\S.\,164.

\gesperrt{Basis of a quadratic field.}
\end{center}

We now wish to apply the results of the general theory of algebraic numbers to the particular case of the fields of degree $2$,
which are also called quadratic fields. A quadratic field $\Omega$
springs from a quadratic equation with integral coefficients, and is
thus derived from the field of the rational numbers by the adjunction
of a square root $\sqrt{d}$. Here $d$ may be assumed to be a positive
or negative integer without square divisor, and the single value
$d = 1$ is to be excluded. According as $d$ is positive or negative,
we obtain a real or an imaginary quadratic field.

We obtain the conjugate field if in all the numbers in $\Omega$ we
change $\sqrt{d}$ into $-\sqrt{d}$. By this, however, we obtain no new
field. The quadratic field is therefore an (absolute) normal field
(\S.\,160).

The numbers of the field $\Omega$ are none other than the quadratic
irrational numbers defined in \S.\,122 of the first Band, which are
all contained in the form $x + y\sqrt{d}$, where $x$ and $y$ denote
integral or fractional rational numbers. Every irrational number
$\omega$ of this field satisfies a quadratic equation

\weq{1}{c\,\omega^2 = a + b\,\omega,}

in which $a, b, c$ are integral rational numbers without common
divisor, of which $c$ may be assumed to be positive, and the
combination

\weq{2}{D = b^2 + 4\,a\,c}
%CONT>>
